% JADE research note: how the system is evaluated and backtested, with the results.
% Build: tectonic docs/research/jade-research.tex  (writes jade-research.pdf beside it)
% Numbers: seed 1337 played year (scripts/step.py --players; tpa/evaluate.py), the make evals run of
% 2026-10-24 (frontend/src/api/fixtures/evals.json), and scripts/backtest.py (tpa/backtest.py, tpa/filings.py).
\documentclass[10pt,letterpaper]{article}
\usepackage[margin=0.9in]{geometry}
\usepackage{newtxtext,newtxmath}
\usepackage{booktabs,multirow,microtype}
\usepackage{url}\usepackage[hidelinks]{hyperref}
\Urlmuskip=0mu plus 1mu
\usepackage{titlesec}
\renewcommand{\thesection}{\Roman{section}}\renewcommand{\thesubsection}{\thesection-\Alph{subsection}}
\titleformat{\section}{\normalfont\large\scshape\centering}{\thesection.}{0.5em}{}
\titleformat{\subsection}{\normalfont\itshape}{\Alph{subsection}.}{0.5em}{}
\renewcommand{\arraystretch}{1.12}
\newcommand{\tabcap}[2]{\begin{center}\small TABLE #1\\[1pt]\textsc{#2}\end{center}\vspace{-6pt}}
\newenvironment{tbl}{\par\medskip\noindent\begin{minipage}{\linewidth}}{\end{minipage}\par\medskip}
\newcommand{\note}[1]{\par\vspace{2pt}\noindent{\footnotesize #1}\par}

\title{\Large JADE: Evaluating an Autonomous Trade Promotion Accountant\\
\large Planted-error evaluation, refusal on real data, and a walk-forward backtest on General Mills' 10-Ks}
\author{JADE team \quad Cal Hacks 13.0, Hyphenate track}
\date{October 2026}

\begin{document}
\maketitle

\begin{abstract}
JADE fills two finance seats at a consumer packaged goods manufacturer: the trade promotion accountant
and the deduction analyst. It accrues promotional obligations on shipment as contra-revenue (ASC 606~\cite{asc606-cpc,asc606-vc}),
matches retailer deductions to programs, clears what the agreement supports, disputes what it does not,
restates closed months when terms change, and packages the close for a person to sign. Every number is
computed by code into an insert-only, bitemporal ledger; a language model only drafts wording. We
evaluate it three ways. \emph{Planted errors:} 42 errors and 8 clean look-alikes hidden in 24 months of
seeded books; JADE catches 42 of 42 with precision 1.000 at month 12, on the development seed and on two
seeds it never saw. \emph{Refusal on real data:} 4,996 real emails and 102 real commercial contracts; it
acts on none. \emph{Real books:} a walk-forward backtest of General Mills' accrued trade promotion
liability, FY2021--FY2026, where the best estimator misses by 9.4\% and the close's analytical review
raises no false alarm on audited figures.
\end{abstract}

\section{Problem}
Trade spend is the second-largest line on a consumer packaged goods income statement after cost of goods,
about 20\% of gross sales~\cite{strategyand}.
General Mills sold \$18.42B in fiscal 2026, 22\% of it to Walmart, and carried \$493M of accrued trade and
coupon promotion liabilities at year end, which its 10-K names as a critical accounting estimate~\cite{gis26}. The
manufacturer books what it owes when the truck leaves and learns what it really owed months later, when
the retailer deducts it from a payment, often at a different rate and with backup that does not tie.
When the estimate trues up, the quarter moves: trade expense timing took 13 points off Q4 adjusted
operating profit growth in FY2025 and added 7 points in FY2026~\cite{gis-q4-25,gis-q4-26}.

Deduction-level data is never public, on either side of the trade. That shapes the evaluation: the
transaction-level behaviour is tested on seeded books with an answer key (Sections~\ref{sec:planted}
and~\ref{sec:learning}), refusal is tested on real public text (Section~\ref{sec:escalating}), and the one
estimate that is public, the year-end liability, is backtested on the real filings (Section~\ref{sec:backtest}).

\section{System under test}
Each step runs one month and handles each source row on the day it lands: read agreements into terms
(each term cites its clause), accrue on shipment, match each deduction to a program, its shipments and
its claim backup, run 12 controls mined from last year's decisions, then clear, dispute or ask. What the
agreement cannot settle becomes a card for the person who owns the decision, with a recommendation; that
person's answer becomes a precedent, reused for the same situation up to the amount approved.
Five verifiability gates (journals balance, reversals mirror, no stale accruals, claims inside their
windows, nothing cleared beyond what was accrued) block a period rather than warn. Segregation of duties is
a capability check that throws: the person who negotiated a deal cannot approve its money. Settlement follows
the trade allowance and deduction workflow of commercial ERPs~\cite{msft} and keeps credit-loss exposure
(ASC 326~\cite{asc326}) apart from promotional liabilities.

\section{Evaluation method}
\subsection{Seeded books and the answer key}
A deterministic generator (seed 1337) produces 24 months of shipments, invoices, returns, deductions,
claim backup, agreements and amendments as documents, and last year's human decisions, for General Mills'
programs at Walmart, Costco and Kroger. Into these it plants 42 errors of known types (over-claims,
duplicates, missing backup, stale accruals, an unentered amendment, self-approval attempts and others) and
8 clean look-alikes that must pass. The answer key (\texttt{data/answer\_key.json}) is never visible to
the system. Volumes are calibrated to the 10-K: every close checks the simulated liability against
General Mills' reported ratio to net sales (2.16\%--3.37\% simulated, against 2.41\% and 2.68\% reported).

\subsection{Scoring}
\texttt{tpa/evaluate.py} scores findings against the key by evidence: a finding counts against a planted
row when its evidence names one of the row's specific ids (deduction, shipment, event, amendment or accrual
key); program ids are too broad to count. A gate failure counts when the period's gate results carry a
violation keyed to the row. We report, cumulatively by month:
\begin{itemize}\itemsep0pt
\item \emph{precision}: findings that match a planted error, over all findings;
\item \emph{recall}: planted errors caught, over planted errors due by that month;
\item \emph{escalation rate}: findings in the month that went to a person;
\item \emph{autonomy}: customer-periods signed with no escalation raised in them.
\end{itemize}
The humans who answer cards in an unattended run are played: they take the agent's recommendation, or a
model plays the role and must pick one of the card's options. Findings on rows injected live during a
demo are scored apart, since the key never planted them.

\subsection{Generalisation and reliability}
The same generator with seeds 42 and 2024, which were never developed against, places the same error types
differently. Holdouts are sized to the claim~\cite{kapoor}: this shows the system is not fitted to one answer key; it does not show it would catch an error
type never planted. Reliability is measured as pass\textasciicircum{}$k$~\cite{taubench}, the chance that all $k$ tries succeed: the full year is run $k=3$ times
from the same seed and the ledger and decision hashes must be identical.

\subsection{Refusal on real data}
\label{sec:escalating-method}
An autonomous accountant's worst failure is acting on input it should not trust. We run real public text
through the two readers that turn documents into source rows: 3,000 business emails
(\texttt{corbt/enron-emails}~\cite{hf-enron}) and 1,996 spam and ham emails (\texttt{SetFit/enron\_spam}~\cite{hf-spam}) through the mail
reader, and 102 commercial contracts from CUAD~\cite{cuad} (\texttt{chenghao/cuad\_qa}~\cite{hf-cuad}), wrapped as amendments, through
the agreement reader. The pass condition is zero actions and zero extracted terms. We also reworded 28
deduction notices by hand to test that the mail reader holds what it cannot parse instead of guessing.

\section{Backtest method}
\label{sec:backtest-method}
The one trade promotion number that is public is the year-end liability. \texttt{tpa/backtest.py} walks
forward one fiscal year-end close at a time over FY2021--FY2026, using figures quoted from General Mills'
10-K filings on SEC EDGAR, CIK 0000040704~\cite{gis21,gis22,gis23,gis24,gis25,gis26} (\texttt{tpa/filings.py}). At each close an estimator sees only the 10-Ks filed
before that year's own 10-K, plus the year's net sales, which the company knows at its close; the estimate
is then scored against the liability the 10-K reported. Four estimators:
\begin{itemize}\itemsep0pt
\item \emph{carry forward}: last year's liability;
\item \emph{annual ratio}: last year's liability over net sales, times this year's net sales (the measure the
close's analytical review uses);
\item \emph{Q4 ratio}: the same on fourth-quarter sales, since claims settle within months;
\item \emph{mean Q4 ratio}: the Q4 ratio averaged over every year public so far.
\end{itemize}
Two of the close's controls are replayed on the real books. Because those books were audited, any flag is a
false alarm: the \emph{analytical review} (the liability ratio against the prior 10-K's, at the policy
tolerance of $\pm$50\%), and a \emph{restatement check} (a year end stated differently a year later).
This scores the top-down estimate and the controls at company level, never deduction matching.

\section{Results}
\subsection{Planted errors}
\label{sec:planted}
\begin{tbl}
\tabcap{I}{Self-improving: planted errors, scored monthly (seed 1337)}
\begin{center}\begin{tabular}{lrrrrrr}
\toprule
Month & Planted & Caught & Precision & Recall & Escalation & Autonomy \\
\midrule
1  & 4  & 4  & 1.000 & 1.000 & 0.400 & 0.667 \\
3  & 10 & 10 & 1.000 & 1.000 & 0.250 & 0.667 \\
5  & 21 & 21 & 1.000 & 1.000 & 0.333 & 0.600 \\
6  & 27 & 27 & 1.000 & 1.000 & 0.143 & 0.611 \\
7  & 30 & 30 & \textbf{0.857}$^\dagger$ & 1.000 & 0.222 & 0.571 \\
8  & 35 & 35 & 1.000 & 1.000 & 0.133 & 0.542 \\
9  & 37 & 37 & 1.000 & 1.000 & 0.800 & 0.518 \\
10 & 40 & 40 & 1.000 & 1.000 & 0.143 & 0.533 \\
12 & 42 & 42 & 1.000 & 1.000 & 0.500 & 0.556 \\
\bottomrule
\end{tabular}\end{center}
\note{Cumulative. $^\dagger$Six July rate conflicts were raised before the planted amendment they trace to
entered the scoring window. Autopilot without the judgment engine catches 40 of 42.}
\end{tbl}

\begin{tbl}
\tabcap{II}{Generalisation and reliability}
\begin{center}\begin{tabular}{lrrr}
\toprule
Run & Caught & Precision & Recall \\
\midrule
Seed 1337 (development) & 42 / 42 & 1.000 & 1.000 \\
Seed 42 (unseen) & 42 / 42 & 1.000 & 1.000 \\
Seed 2024 (unseen) & 42 / 42 & 1.000 & 1.000 \\
\bottomrule
\end{tabular}\end{center}
\note{Reliability, pass\textasciicircum3 on seed 1337: 3 of 3 full runs identical (979 ledger lines, 75 findings,
ledger hash \texttt{fa85730190eb4cc1}, decisions hash \texttt{9b39dd7e5a9f7ba3}).}
\end{tbl}

The escalation rate does not fall steadily because new kinds of situations keep arriving through the
year (a learned control in September, a new ambiguous clause); repeated situations do not escalate,
as Table~III shows.

\subsection{Self-learning}
\label{sec:learning}
\begin{tbl}
\tabcap{III}{Self-learning: answers reused as precedent}
\begin{center}\begin{tabular}{lrr}
\toprule
Situation & Cards & Share \\
\midrule
First-time question, asked a person & 17 & --- \\
Repeat situation & 40 & 100.0\% \\
\quad settled silently by precedent & \textbf{37} & \textbf{92.5\%} \\
\quad asked again & 3 & 7.5\% \\
\midrule
\multicolumn{3}{l}{\emph{Why the three were asked again}} \\
FN-0061 \; \$42,413.95 vs.\ approved \$41,319.56 & \multicolumn{2}{r}{over approval} \\
FN-0072 \; \$48,166.44 vs.\ approved \$42,413.95 & \multicolumn{2}{r}{over approval} \\
FN-0075 \; \$58,315.83 vs.\ approved \$48,166.44 & \multicolumn{2}{r}{over approval} \\
\bottomrule
\end{tabular}\end{center}
\note{Seed 1337, twelve played months. A precedent covers the same situation up to the amount a person
approved; a larger exposure goes back to a person.}
\end{tbl}

\subsection{Self-escalating: refusal on real data}
\label{sec:escalating}
\begin{tbl}
\tabcap{IV}{Self-escalating: inputs it must not act on}
\begin{center}\begin{tabular}{llrrr}
\toprule
Dataset & Source & $n$ & Acted on & Held / refused \\
\midrule
\texttt{corbt/enron-emails} & Hugging Face & 3,000 & \textbf{0} & 3,000 \\
\texttt{SetFit/enron\_spam} & Hugging Face & 1,996 & \textbf{0} & 1,996 \\
CUAD contracts (\texttt{chenghao/cuad\_qa}) & Hugging Face & 102 & 0 terms & 102 \\
\midrule
Reworded deduction notices & ours & 28 & \multicolumn{2}{r}{25 parsed $\rightarrow$ 28 after 2 fixes} \\
\bottomrule
\end{tabular}\end{center}
\note{Rewording found two real parser gaps (bulleted labels, HTML-escaped senders), fixed in
\texttt{tpa/mail.py}.}
\end{tbl}

\subsection{Real books: the 10-K backtest}
\label{sec:backtest}
\begin{tbl}
\tabcap{V}{Estimate at each close against the reported liability (\$M)}
\begin{center}\begin{tabular}{lrrrrrr}
\toprule
FY & Reported & Carry fwd & Annual ratio & Q4 ratio & Mean Q4 & Review drift \\
\midrule
2021 & 508 & 471.0 & 484.4 & 424.2 & 424.2 & $+4.9\%$ \\
2022 & 420 & 508.0 & 532.3 & 549.3 & 504.0 & $-21.1\%$ \\
2023 & 394 & 420.0 & 444.4 & 431.9 & 489.5 & $-11.3\%$ \\
2024 & 425 & 394.0 & 389.4 & 369.2 & 436.4 & $+9.2\%$ \\
2025 & 470 & 425.0 & 417.1 & 410.8 & 419.6 & $+12.7\%$ \\
2026 & 493 & 470.0 & 444.4 & 475.5 & 433.0 & $+10.9\%$ \\
\bottomrule
\end{tabular}\end{center}
\note{Each estimate made from filings public before that year's 10-K. Review drift is the year's
liability-to-sales ratio against the prior year's.}
\end{tbl}

\begin{tbl}
\tabcap{VI}{Real books: General Mills 10-K backtest, summary}
\begin{center}\begin{tabular}{lrrrr}
\toprule
Estimate at each close & MAE (\$M) & Mean error & Bias & Worst year \\
\midrule
Carry last year forward & \textbf{41.7} & \textbf{9.4\%} & $-0.2\%$ & 21.0\% \\
Annual ratio to net sales & 53.9 & 12.3\% & $+0.9\%$ & 26.7\% \\
Q4 ratio to net sales & 63.9 & 14.4\% & $-0.9\%$ & 30.8\% \\
Mean Q4 ratio & 64.2 & 14.4\% & $+1.2\%$ & 24.2\% \\
\bottomrule
\end{tabular}\end{center}
\note{Accrued trade and coupon promotions at each fiscal year end, FY2021--FY2026 (six closes,
\$394M--\$508M). Analytical review at $\pm$50\%: 0 false alarms; largest drift $-21.1\%$ (FY2022).
Restatements of a reported year end: none.}
\end{tbl}

Two findings follow. First, the audited liability swings by up to 21.1\% a year relative to sales, so a
$\pm$50\% analytical review can never fail on these books; a tolerance near $\pm$22\% would be calibrated
on the filings. Second, no top-down rule beats carrying last year forward, and scaling by sales does worse:
the balance follows which programs run and how late retailers claim, which is why JADE accrues program by
program rather than as a share of sales.

\section{Limitations}
Transaction-level results are on simulated books: the generator and the error types are ours, and
unseen seeds vary placement, not kind. The played humans accept recommendations; real reviewers will
disagree, and the precedent store has only been exercised on the situations it saw. The backtest scores
six company-level year ends, not deduction matching, and FY2021's estimate rests on one prior filing.
The tolerances used are company policy settings, not GAAP requirements.

\section{Reproducing}
\begin{flushleft}\small
\texttt{make db-up \&\& make db-scratch \&\& uv sync \&\& make seed} \quad 24 months, seed 1337\\
\texttt{make players} \quad the played year and its scores (Tables I and III)\\
\texttt{uv run python scripts/evaluate.py --as-of month-1 --as-of month-9} \quad scoring at chosen months\\
\texttt{make backtest} \quad Tables V and VI, no database needed\\
\texttt{tectonic docs/research/jade-research.tex} \quad this paper
\end{flushleft}

\section*{References}
\begingroup\small
\renewcommand{\refname}{}\vspace{-24pt}
\begin{thebibliography}{99}\setlength{\itemsep}{1pt}
\bibitem{gis26} General Mills, Inc., Form 10-K for the fiscal year ended May 31, 2026, filed 2026-07-01, SEC accession 0001628280-26-046466. \url{https://www.sec.gov/Archives/edgar/data/40704/000162828026046466/gis-20260531.htm}
\bibitem{gis25} General Mills, Inc., Form 10-K, fiscal 2025, filed 2025-06-26, accession 0001193125-25-147079. \url{https://www.sec.gov/Archives/edgar/data/40704/000119312525147079/d938443d10k.htm}
\bibitem{gis24} General Mills, Inc., Form 10-K, fiscal 2024, filed 2024-06-26, accession 0001193125-24-168943. \url{https://www.sec.gov/Archives/edgar/data/40704/000119312524168943/}
\bibitem{gis23} General Mills, Inc., Form 10-K, fiscal 2023, filed 2023-06-28, accession 0001193125-23-177500. \url{https://www.sec.gov/Archives/edgar/data/40704/000119312523177500/}
\bibitem{gis22} General Mills, Inc., Form 10-K, fiscal 2022, filed 2022-06-30, accession 0001193125-22-185257. \url{https://www.sec.gov/Archives/edgar/data/40704/000119312522185257/}
\bibitem{gis21} General Mills, Inc., Form 10-K, fiscal 2021, filed 2021-06-30, accession 0001193125-21-204830; fiscal 2020, filed 2020-07-02, accession 0001193125-20-186469. \url{https://www.sec.gov/Archives/edgar/data/40704/000119312521204830/}
\bibitem{gis-q4-25} General Mills, Inc., Q4 fiscal 2025 earnings release, Form 8-K Ex.~99, 2025-06-25, accession 0001193125-25-146132. \url{https://www.sec.gov/Archives/edgar/data/40704/000119312525146132/}
\bibitem{gis-q4-26} General Mills, Inc., Q4 fiscal 2026 earnings release, Form 8-K Ex.~99, 2026-07-01, accession 0001628280-26-046337. \url{https://www.sec.gov/Archives/edgar/data/40704/000162828026046337/}
\bibitem{strategyand} Strategy\& (PwC), ``Zero-based trade for CPG leaders,'' 2017. \url{https://www.strategyand.pwc.com/gx/en/insights/2017/zero-based-trade-for-cpg-leaders/zero-based-trade-for-cpg-leaders.pdf}
\bibitem{asc606-cpc} Deloitte, \emph{Roadmap: Revenue Recognition}, \S6.6, Consideration payable to a customer (ASC 606-10-32-25 to 32-27). \url{https://dart.deloitte.com/USDART/home/codification/revenue/asc606-10/roadmap-revenue-recognition/chapter-6-step-3-determine-transaction/6-6-consideration-payable-a-customer}
\bibitem{asc606-vc} Deloitte, \emph{Roadmap: Revenue Recognition}, \S6.3, Variable consideration (ASC 606-10-32-5 to 32-13). \url{https://dart.deloitte.com/USDART/home/codification/revenue/asc606-10/roadmap-revenue-recognition/chapter-6-step-3-determine-transaction/6-3-variable-consideration}
\bibitem{asc326} Deloitte, \emph{Roadmap: Current Expected Credit Losses}, \S5.2, Trade receivables and contract assets (ASC 326). \url{https://dart.deloitte.com/USDART/home/codification/assets/32x/asc326-10/roadmap-credit-losses-cecl/chapter-5-application-cecl-model-off/5-2-trade-receivables-contract-assets}
\bibitem{msft} Microsoft, ``Trade allowance management and deduction settlement,'' Dynamics 365 Supply Chain Management documentation. \url{https://learn.microsoft.com/en-us/dynamics365/supply-chain/sales-marketing/trade-allowance}
\bibitem{taubench} S.~Yao, N.~Shinn, P.~Razavi, K.~Narasimhan, ``$\tau$-bench: A benchmark for tool-agent-user interaction in real-world domains,'' arXiv:2406.12045, 2024.
\bibitem{kapoor} S.~Kapoor, B.~Stroebl, Z.~S.~Siegel, N.~Nadgir, A.~Narayanan, ``AI agents that matter,'' arXiv:2407.01502, 2024.
\bibitem{cuad} D.~Hendrycks, C.~Burns, A.~Chen, S.~Ball, ``CUAD: An expert-annotated NLP dataset for legal contract review,'' NeurIPS Datasets and Benchmarks, arXiv:2103.06268, 2021.
\bibitem{hf-cuad} \texttt{chenghao/cuad\_qa}, Hugging Face dataset. \url{https://huggingface.co/datasets/chenghao/cuad_qa}
\bibitem{hf-enron} \texttt{corbt/enron-emails}, Hugging Face dataset (the Enron email corpus). \url{https://huggingface.co/datasets/corbt/enron-emails}
\bibitem{hf-spam} \texttt{SetFit/enron\_spam}, Hugging Face dataset. \url{https://huggingface.co/datasets/SetFit/enron_spam}
\end{thebibliography}
\endgroup

\end{document}
